% ======================================================================
%  PART 8 -- Forecasting
% ======================================================================
\section{Forecasting}
\label{sec:forecasting}

The purpose of time series analysis is often to \emph{forecast} future values
from the data observed up to now. The recipe has three steps: (1) assume the
form of the model; (2) estimate its parameters; (3) use the assumed structure
and the estimated parameters to form predictions. In this part we take the
model structure and the parameters as \textbf{known} and predict $h$ lags
ahead, i.e.\ we want $X_{n+h}$ from $X_1,\dots,X_n$. The integer $h>0$ is the
\textbf{lead} (or \textbf{horizon}) and $n$ is the \textbf{origin}. The central
theme is that, for the second-order theory we use, the optimal predictor depends
\emph{only} on the covariance structure $\{\acvs_\tau\}$, and for causal ARMA
models it is computed cleanly by setting future innovations to zero.

% ----------------------------------------------------------------------
\subsection{Definitions}

\begin{definition}{Prediction mean square error}{p8-pmse}
Let $g(X_1,\dots,X_n)$ be any estimator (predictor) of $X_{n+h}$ based on the
observations $X_1,\dots,X_n$. Its performance is measured by the
\textbf{prediction mean square error}
\[
P^{\,n}_{n+h}\;:=\;\E\!\left[\big(X_{n+h}-g(X_1,\dots,X_n)\big)^2\right].
\]
Smaller $P^{\,n}_{n+h}$ means a better forecast; minimising it over a class of
predictors defines the ``best'' predictor in that class.
\end{definition}

\begin{definition}{Best linear predictor}{p8-blp}
The \textbf{best linear predictor} of $X_{n+h}$ based on $X_1,\dots,X_n$ is the
linear function
\[
X^{\,n}_{n+h}\;=\;P_{X_1,\dots,X_n}(X_{n+h})\;=\;\sum_{j=1}^{n}\beta_j X_j
\]
that minimises the prediction mean square error among all linear combinations of
$X_1,\dots,X_n$. We write $P_{X_1,\dots,X_n}(\cdot)$ for the projection onto the
linear span of the data.
\begin{itemize}
  \item For \textbf{Gaussian} time series the conditional expectation
        $\E[X_{n+h}\mid X_1,\dots,X_n]$ is itself linear in the data, so the best
        linear predictor coincides with the (overall) best predictor.
  \item In general the conditional expectation is a \emph{nonlinear} function of
        the data; for tractability one restricts attention to linear predictors,
        which need only second-order information.
\end{itemize}
\end{definition}

\begin{definition}{$h$-step forecast error and its variance}{p8-ferr}
For the best linear predictor $X^{\,n}_{n+h}$ define the \textbf{$h$-step ahead
forecast error}
\[
e_n(h)\;:=\;X_{n+h}-X^{\,n}_{n+h}.
\]
It is unbiased, $\E[e_n(h)]=0$, and its variance equals the prediction MSE,
$\Var\big(e_n(h)\big)=P^{\,n}_{n+h}$. For a causal ARMA process with
$MA(\infty)$ weights $\{\psi_j\}$ one has
$e_n(h)=\sum_{j=0}^{h-1}\psi_j\eps_{n+h-j}$ and
$\Var\big(e_n(h)\big)=\sigma^2\sum_{j=0}^{h-1}\psi_j^2$.
\end{definition}

\begin{definition}{Sources of prediction uncertainty}{p8-srcunc}
There are \textbf{three} components of uncertainty in a forecast:
\begin{enumerate}
  \item \textbf{model uncertainty} --- the assumed model form may be wrong;
  \item \textbf{estimation uncertainty} --- the parameters are estimated, not known;
  \item \textbf{innovation uncertainty} --- the inherent randomness of the future
        innovations $\eps_{n+1},\dots,\eps_{n+h}$.
\end{enumerate}
Throughout we assume the model is \emph{known} and the parameters are estimated
\emph{without error}, which removes the first two sources and leaves only
innovation uncertainty. (Bootstrap methods can incorporate the first two, but
are not covered here.)
\end{definition}

\begin{definition}{Gaussian prediction interval}{p8-pint}
For a \textbf{Gaussian} process the $h$-step error $e_n(h)$ is normal with mean
$0$ and variance $\Var(e_n(h))$, so a $(1-\alpha)$ \textbf{prediction interval}
for $X_{n+h}$ based on $X_1,\dots,X_n$ is
\[
X^{\,n}_{n+h}\;\pm\;z_{1-\alpha/2}\,\sqrt{\Var\big(e_n(h)\big)},
\]
where $z_{1-\alpha/2}$ is the $(1-\alpha/2)$ quantile of the standard normal
distribution (e.g.\ $z_{0.975}\approx1.96$ for a $95\%$ interval).
\end{definition}

% ----------------------------------------------------------------------
\subsection{Theorems \& Propositions}

\begin{lemma}{Conditional expectation minimises the prediction MSE}{p8-condexp}
The conditional expectation
\[
g(X_1,\dots,X_n)=\E\big[X_{n+h}\mid X_1,\dots,X_n\big]
\]
minimises the prediction mean square error for $X_{n+h}$ based on
$X_1,\dots,X_n$.\\[3pt]
\textit{Proof idea:} For any real random variable $Y$ with $\mu=\E[Y]$ and any
constant $c$,
\[
\MSE(c)=\E\big[(Y-c)^2\big]=\E\big[(Y-\mu+\mu-c)^2\big]
=\E\big[(Y-\mu)^2\big]+(\mu-c)^2=\Var(Y)+(\mu-c)^2,
\]
which is minimised at $c=\mu=\E[Y]$ (the cross term vanishes because
$\E[Y-\mu]=0$). Applying this \emph{conditionally on} $X_1,\dots,X_n$ and using
the tower property
$\E[(X_{n+h}-g)^2]=\E\big[\E[(X_{n+h}-g)^2\mid X_1,\dots,X_n]\big]$, the inner
MSE is minimised pointwise by taking $g=\E[X_{n+h}\mid X_1,\dots,X_n]$.\\[2pt]
\textit{Why:} this is the theoretical gold standard for forecasting. For
Gaussian processes the conditional expectation is linear, so the best predictor
\emph{is} the best linear predictor; in general it motivates restricting to
linear predictors, which only need the ACVS.
\end{lemma}

\begin{theorem}{Prediction equations (zero-mean stationary process)}{p8-predeq}
For a zero-mean stationary process $\{X_t\}$, the best linear predictor
$X^{\,n}_{n+h}=\sum_{j=1}^n\beta_j X_j$ is found by solving for
$\beta_1,\dots,\beta_n$ the \textbf{prediction equations}
\[
\E\!\big[(X_{n+h}-X^{\,n}_{n+h})\,X_k\big]=0,\qquad k=1,\dots,n,
\]
i.e.\ the forecast error is orthogonal to every observation. Equivalently, in
matrix form $\Gamma_n\beta=\acvs_{[h]}$:
\[
\begin{pmatrix}
\acvs_0 & \acvs_1 & \cdots & \acvs_{n-1}\\
\acvs_1 & \acvs_0 & \cdots & \acvs_{n-2}\\
\vdots & \vdots & \ddots & \vdots\\
\acvs_{n-1} & \acvs_{n-2} & \cdots & \acvs_0
\end{pmatrix}
\begin{pmatrix}\beta_1\\ \beta_2\\ \vdots\\ \beta_n\end{pmatrix}
=\begin{pmatrix}\acvs_{n+h-1}\\ \acvs_{n+h-2}\\ \vdots\\ \acvs_{h}\end{pmatrix}.
\]
\textit{Proof idea:} the best linear predictor is the orthogonal projection of
$X_{n+h}$ onto $\mathrm{span}\{X_1,\dots,X_n\}$; the residual
$X_{n+h}-X^{\,n}_{n+h}$ must be orthogonal (uncorrelated) to each $X_k$.
Expanding $\E[(X_{n+h}-\sum_j\beta_jX_j)X_k]=0$ gives
$\acvs_{n+h-k}=\sum_{j=1}^n\beta_j\acvs_{j-k}$, which is row $k$ of
$\Gamma_n\beta=\acvs_{[h]}$.\\[2pt]
\textit{Why:} this is the master equation for linear prediction of \emph{any}
stationary process --- it shows the optimal forecast uses only the ACVS.
\end{theorem}

\begin{corollary}{Matrix solution and prediction MSE}{p8-matsol}
If $\Gamma_n$ is invertible, then with $X=(X_1,\dots,X_n)\Tr$ the best linear
predictor and its MSE are
\[
X^{\,n}_{n+h}=\acvs_{[h]}\Tr\,\Gamma_n^{-1}\,X,
\qquad
P^{\,n}_{n+h}=\acvs_0-\acvs_{[h]}\Tr\,\Gamma_n^{-1}\,\acvs_{[h]} .
\]
\textit{Proof idea:} solve $\Gamma_n\beta=\acvs_{[h]}$ for
$\beta=\Gamma_n^{-1}\acvs_{[h]}$ and substitute into
$X^{\,n}_{n+h}=\beta\Tr X$ and into
$P^{\,n}_{n+h}=\acvs_0-2\beta\Tr\acvs_{[h]}+\beta\Tr\Gamma_n\beta$; the last
two terms collapse using $\Gamma_n\beta=\acvs_{[h]}$.\\[2pt]
\textit{Why:} gives the closed-form forecast for short series.
\begin{itemize}
  \item $\Gamma_n$ is invertible for \emph{invertible} ARMA models (and not in
        all other cases), but the best prediction $X^{\,n}_{n+h}$ is
        \textbf{always unique} even when $\Gamma_n$ is singular.
  \item Computing $\acvs_{[h]}\Tr\Gamma_n^{-1}X$ is inefficient for large $n$
        because the $n\times n$ matrix must be inverted; the Durbin--Levinson
        and innovations algorithms avoid this (Shumway \& Stoffer, \S3.5).
\end{itemize}
\end{corollary}

\begin{theorem}{Prediction equations with non-zero mean (and intercept $\beta_0$)}{p8-predeqmean}
For a stationary process with $\mu=\E[X_t]$, write
$X^{\,n}_{n+h}=\beta_0+\sum_{j=1}^n\beta_j X_j$. The coefficients solve the
prediction equations $\E[(X_{n+h}-X^{\,n}_{n+h})X_k]=0$ for $k=0,1,\dots,n$
(with the convention $X_0\equiv1$), which give
\[
\Gamma_n\beta=\acvs_{[h]},
\qquad
\beta_0=\mu\big(1-\beta\Tr\mathbf 1_n\big),
\]
and, when $\Gamma_n$ is invertible,
\[
X^{\,n}_{n+h}=\mu+\acvs_{[h]}\Tr\,\Gamma_n^{-1}(X-\mu\mathbf 1_n),
\qquad
P^{\,n}_{n+h}=\acvs_0-\acvs_{[h]}\Tr\,\Gamma_n^{-1}\,\acvs_{[h]} .
\]
\textit{Proof idea:} the predictor has the form
$X^{\,n}_{n+h}=\mu+\beta\Tr(X-\mu\mathbf 1_n)$, so centring removes the mean and
reduces the problem to the zero-mean case --- there is no loss of generality in
assuming $\mu=0$, where $\beta_0=0$. The MSE is \emph{identical} to the
zero-mean case because the constant shift does not affect the error variance.\\[2pt]
\textit{Why:} real data are rarely centred; this is the version actually used in
practice.
\end{theorem}

\begin{theorem}{Best linear predictor of a causal ARMA via $\psi$-weights}{p8-armapred}
For a causal ARMA process with $MA(\infty)$ representation
$X_t=\sum_{j=0}^\infty\psi_j\eps_{t-j}$, the best linear predictor of $X_{n+h}$
based on $X_1,\dots,X_n$ is
\[
X^{\,n}_{n+h}=\sum_{j=h}^{\infty}\psi_j\,\eps_{n+h-j}
=\psi_h\eps_n+\psi_{h+1}\eps_{n-1}+\cdots,
\]
with prediction mean square error
\[
P^{\,n}_{n+h}=\sigma^2\sum_{j=0}^{h-1}\psi_j^2 .
\]
\textit{Proof idea:} write the predictor as a linear filter of the past
innovations, $X^{\,n}_{n+h}=\sum_{j=0}^\infty c_j\eps_{n-j}$. Then
\[
\E\big[(X_{n+h}-X^{\,n}_{n+h})^2\big]
=\sigma^2\Big(\sum_{j=0}^{h-1}\psi_j^2+\sum_{j=h}^\infty(\psi_j-c_{j-h})^2\Big),
\]
which is minimised by matching $c_{j}=\psi_{j+h}$, killing the second sum and
leaving $\sigma^2\sum_{j=0}^{h-1}\psi_j^2$. Practically:
\textbf{set the unrealised future innovations $\eps_{n+1},\dots,\eps_{n+h}$ to
zero} in $X_{n+h}=\sum_{j=0}^\infty\psi_j\eps_{n+h-j}$.\\[2pt]
\textit{Why:} this is the engine of all ARMA forecasting. It also shows that for
a causal ARMA, $\psi_j\to0$ exponentially so $X^{\,n}_{n+h}\to\mu$, and
$P^{\,n}_{n+h}=\sigma^2\sum_{j=0}^{h-1}\psi_j^2\uparrow\sigma^2\sum_{j=0}^\infty\psi_j^2=\acvs_0$
as $h\to\infty$.
\end{theorem}

\begin{proposition}{Limiting behaviour of ARMA forecasts}{p8-armalimit}
For a causal ARMA model with mean $\mu$,
$X^{\,n}_{n+h}=\mu+\sum_{j=h}^\infty\psi_j\eps_{n+h-j}$ with $\psi_j\to0$
exponentially fast, hence
\[
X^{\,n}_{n+h}\longrightarrow\mu,
\qquad
\Var\big(e_n(h)\big)=\sigma^2\sum_{j=0}^{h-1}\psi_j^2\longrightarrow\acvs_0,
\qquad h\to\infty .
\]
\textit{Proof idea:} immediate from the $\psi$-weight theorem and exponential
decay of $\psi_j$ (a consequence of causality, roots of $\Phi$ outside the unit
circle).\\[2pt]
\textit{Why:} quantifies ``the distant future is the unconditional
distribution'' --- the point forecast converges to the mean and the error
variance converges to the marginal variance $\acvs_0$.
\end{proposition}

\begin{proposition}{Forecasts for integrated (ARIMA) processes}{p8-arimapred}
Let $\{X_t\}$ be ARIMA$(p,d,q)$, $\Phi(B)(1-B)^dX_t=\Theta(B)\eps_t$. Then
$Z_t=\diff^{\,d}X_t$ is the ARMA$(p,q)$ process
$Z_t=\sum_{j=1}^p\phi_jZ_{t-j}+\eps_t-\sum_{j=1}^q\theta_j\eps_{t-j}$, and one
forecasts $X$ by undoing the differencing,
\[
\diff^{\,d}X^{\,n}_{n+h}=Z^{\,n}_{n+h}.
\]
For $d=1$ (so $X_{n+h}-X_{n+h-1}=Z_{n+h}$),
\[
X^{\,n}_{n+h}=X_n+\sum_{j=0}^{h-1}Z^{\,n}_{n+h-j},\qquad h\ge1 .
\]
If $\{Z_t\}$ has mean $\alpha\ne0$ then $Z^{\,n}_{n+h}\to\alpha$ and
$X^{\,n}_{n+h}\approx X_n+\alpha h$ for large $h$ (a deterministic drift line).\\[2pt]
\textit{Proof idea:} apply the linear projection operator $P_{X_1,\dots,X_n}$,
which is linear, to $\diff^{\,d}X_{n+h}=Z_{n+h}$; since differencing is a fixed
linear filter it commutes with the projection, giving
$\diff^{\,d}X^{\,n}_{n+h}=Z^{\,n}_{n+h}$. Telescoping the $d=1$ identity
$X_{n+h}=X_{n+h-1}+Z_{n+h}$ down to the last observation $X_n$ yields
$X^{\,n}_{n+h}=X_n+\sum_{j=0}^{h-1}Z^{\,n}_{n+h-j}$.\\[2pt]
\textit{Why:} reduces ARIMA forecasting to ARMA forecasting on the differenced
series, then ``re-integrates''.
\begin{itemize}
  \item For ARIMA models the prediction MSE \textbf{diverges}:
        $P^{\,n}_{n+h}\to\infty$ as $h\to\infty$ (no mean reversion --- the
        process is non-stationary).
  \item For SARIMA models with $d=D=1$ and $\alpha=0$ the forecast is linear and
        seasonal in $h$.
\end{itemize}
\end{proposition}

% ----------------------------------------------------------------------
% ----------------------------------------------------------------------
